% !PRINTED 131
\weberpage{131}

Wenn wir die Ungleichungen (13) so schreiben
\[
\begin{aligned}
 a-a'&>\frac{a^2}{4Q},\\
 a'-a''&>\frac{a'^2}{4Q},\\
 &\hspace{1.5em}\vdots\\
 a^{(\nu-1)}-a^{(\nu)}&>\frac{(a^{(\nu-1)})^2}{4Q},
\end{aligned}
\]
und sie dann von der $(\mu+1)$ten bis zur $(\nu+1)$ten addiren, so folgt
\[
 (a^{(\mu)})^2+(a^{(\mu+1)})^2+\cdots+(a^{(\nu)})^2
 <4Q\bigl(a^{(\mu)}-a^{(\nu+1)}\bigr),
\]
also, wenn $\mu$ und $\nu$ zwei ins Unendliche wachsende Zahlen bedeuten,
\begin{equation}
 \lim\bigl((a^{(\mu)})^2+(a^{(\mu+1)})^2+\cdots+(a^{(\nu)})^2\bigr)=0.
\tag{16}
\end{equation}

Für die Reihe der Zahlen $\alpha,\alpha',\alpha'',\ldots$ ergiebt sich aus (11) das Bildungsgesetz
\begin{equation}
\begin{aligned}
 \alpha-\alpha'&=\frac{a f(\alpha)}{2Q f'(\alpha)},\\
 \alpha'-\alpha''&=\frac{a' f(\alpha')}{2Q f'(\alpha')},\\
 &\hspace{1.5em}\vdots\\
 \alpha^{(\nu-1)}-\alpha^{(\nu)}&=
 \frac{a^{(\nu-1)} f(\alpha^{(\nu-1)})}{2Q f'(\alpha^{(\nu-1)})}.
\end{aligned}
\tag{17}
\end{equation}
Wenn man eine beliebige Anzahl dieser Gleichungen von der $(\mu+1)$ten bis zur $(\nu+1)$ten addirt, so folgt
\begin{equation}
 \alpha^{(\mu)}-\alpha^{(\nu+1)}=
 \frac{a^{(\mu)}f(\alpha^{(\mu)})}{2Q f'(\alpha^{(\mu)})}+\cdots+
 \frac{a^{(\nu)}f(\alpha^{(\nu)})}{2Q f'(\alpha^{(\nu)})}.
\tag{18}
\end{equation}
Daraus ergiebt sich, da die sämmtlichen
\[
 |f'(\alpha)|>k
\]
sind, für den absoluten Werth dieser Differenz
\begin{equation}
 |\alpha^{(\mu)}-\alpha^{(\nu+1)}|<
 \frac{(a^{(\mu)})^2+(a^{(\mu+1)})^2+\cdots+(a^{(\nu)})^2}{2Qk},
\tag{19}
\end{equation}
also mit Hülfe von (16)
\begin{equation}
 \lim |\alpha^{(\mu)}-\alpha^{(\nu+1)}|=0.
\tag{20}
\end{equation}

Damit ist ausgedrückt (nach dem in der Einleitung erklärten Begriff der Zahlenreihe, S. 15), dass sich die Zahlen $\alpha,\alpha',\alpha'',\ldots$ einer bestimmten Grenze nähern, die wir mit $\xi$ bezeichnen wollen.
%CONT>>
